\section{Background}

Evaluating LLM agents on the actions they take, alongside their stated moral judgments, has become a central concern for AI safety evaluation \cite{pan2023machiavelli, lynch2025agentic}. To situate our study, we review LLMs' escalation tendencies in high-stakes simulations and the relationship between their ethical reasoning and actions.

\subsection{LLM Escalation in High-Stakes Simulations}

Studies of scripted wargames have found LLMs' escalation tendency in nuclear arms races. For example, \citet{rivera2024escalation} observe arms-race dynamics and occasional nuclear use, \citet{lamparth2024human} find more aggressive U.S.-China crisis recommendations than expert baselines, and \citet{payne2026aiarms} reports 95\% nuclear-threshold crossing from SOTA models' self-play. Escalation patterns differ across model families in wider high-stakes simulations \cite{shrivastava2024freeform}, showing distinguishable signatures in moral robustness \cite{costa2026moral} or strategic heuristics \cite{defortuny2025strategic}. Moreover, stronger reasoning capability may not reliably mitigate catastrophic or deceptive behaviors \cite{xu2025nuclear}, or producing desirable collaboration \cite{piedrahita2025corrupted}.

Predefined crisis states and action spaces support controlled comparisons, but they can also foreground nuclear escalation, making it difficult to interpret differences across reported findings. For example, the scenarios in \citet{rivera2024escalation} and \citet{payne2026aiarms} make nuclear escalation a salient action. By contrast, models in \citet{solopova2026strategic}'s real-world geopolitical vignettes (e.g., trade wars and arctic tensions) and those in \citet{lamparth2024human} did not escalate to nuclear use, yet these scenarios lacked nuclear options. Prompt scaffolding and repeated experimentation can flip outcomes. In \citet{elbaum2025managing}'s replication of \citet{rivera2024escalation}, a reflection prompt asking for ``private thoughts about de-escalation strategies to reduce risk'' substantially reduced escalation. In general social-science experiments, repeated runs alone can alter LLMs' decision-making outcomes \cite{zhou2025pimmur}.

\subsection{LLMs' Ethical Reasoning}

In scripted single-turn dilemmas, LLMs can display measurable procedural competence with canonical moral frameworks. MoReBench studies report that LLMs lean toward act utilitarianism and deontology (i.e., rule-based distinctions between right and wrong) \cite{chiu2026morebench, rachels2012elements}. \citet{seror2025moral}'s preference test suggests rational structure in LLMs' moral thinking, with many models exhibiting ``nearly stable moral principles''. In trolley-problem probes, \citet{samway2025consequentialist} show that models' pre-decision chain-of-thought traces skew deontological, while their post-hoc explanations skew consequentialist (i.e., judging an action by its consequences). Models can also respond to ethical instructions by self-correcting \cite{ganguli2023moral, liu2024intrinsic}, activating latent moral concepts that stabilize internal representations \cite{liu2025convergence, lee2026selfcorrection}.

However, perspective shifts, protocol choices, and direction-flipped context can all shift LLMs' elicited moral responses \cite{vannuenen2026fragility, sauter2026rules, blandfort2026moral}. Procedural variations in prisoner's-dilemma settings can produce variability in LLM outputs \cite{robinson2025framing}. In multi-round and accumulated-context settings, LLMs' moral reasoning trajectories are often inconsistent in value preferences \cite{wu2025staircase} and moral frameworks \cite{huang2026moraltrajectories}, and unstable trajectories are more susceptible to persuasive attacks \cite{huang2026moraltrajectories}. Increased capabilities in math or game-strategy domains may not help: effective cognitive patterns in problem-solving often fail to transfer to value reasoning \cite{lee2026clash}.

In agentic settings, LLMs' moral behavior can decouple from verbalized ethical reasoning, reducing prompt-based interventions' effectiveness. When ethics and payoff conflict in decision-making contexts, models may not consistently behave morally \cite{pan2023machiavelli}. In prisoner's-dilemma and public-goods games, survival-oriented manipulations can further depress cooperation \cite{backmann2025ethics}. Activation steering toward altruism can shift LLMs' choices and post-hoc justifications, yet altruistic rhetoric can coexist with unchanged selfish play \cite{sun2026personavectors}. When models are instructed to pursue goals that require harmful action, they verbalize ethical content in reasoning trails while still executing the harm, even under direct instructions to avoid it \cite{lynch2025agentic}. Ethical prompts can also impose costs when incompetence masquerades as adherence \cite{potham2025hierarchical}.
